\section{Security Model and Course Connections}
The attacker controls one project document. It does not control the direct task, initial source, executable tests, installed skills, collector or independent checker. The objective is to induce protected-test changes in place of a source repair. The investigator has native session exports, events, tool arguments and returns, initial hashes, final files and checker results, but no observation of the model's internal decision process.

The Module 1 connection is a trust-boundary error: lower-trust project content is promoted into authority over a user-authorized action. The attacker reuses legitimate editing capability. This is an analogy to injected data steering execution; it is not a claim that tokens execute machine instructions or that a prompt restriction has hardware-enforced semantics.

The Module 2 connection is post-event dependency reconstruction. Tool calls take the place of coarse execution events; file versions and returns are objects; a protected modification is an investigation starting point. Kernel-level provenance can observe system calls, whereas our record observes a whole shell command as one tool call. We therefore label command-derived file relations as inferences and do not claim complete system-call coverage.

The two external controls answer different questions. A protected-file hash comparison checks the final state, while graph edges organize the recorded or inferred actions that could explain it. Neither the agent's statement of success nor the presence of suspicious wording is used alone as a verdict on legitimate task completion.
